% Une seule diapositive de synthèse ; les méthodes citées ne sont pas HGP.
\begin{frame}[t]{Quelques résultats encourageants}
\small
\setlength{\tabcolsep}{6pt}
\renewcommand{\arraystretch}{1.08}
\begin{tabular}{@{}>{\raggedright\arraybackslash}p{.16\textwidth}
                  >{\raggedright\arraybackslash}p{.50\textwidth}
                  >{\raggedright\arraybackslash}p{.27\textwidth}@{}}
\toprule
\textbf{Idée} & \textbf{Résultat publié} & \textbf{Pour notre projet} \\
\midrule
\textbf{Superpoints} &
SPT : segmentation sur KITTI-360 et DALES à partir de régions. \citb{spt23} &
Apprendre sur des régions. \\
\addlinespace[.18cm]
\textbf{Hiérarchie} &
SPT : \textbf{$-5{,}1$ points de mIoU} avec un seul niveau (KITTI-360, validation). \citb{spt23} &
Conserver plusieurs échelles. \\
\addlinespace[.18cm]
\textbf{Distances} &
BPS + petit réseau : précision comparable à PointNet sur ModelNet40. \citb{bps19} &
Tester le code sans grand modèle. \\
\addlinespace[.18cm]
\textbf{Attributs} &
PolyhedronNet : gain avec les attributs de faces (ShapeNet-P). \citb{poly25} &
Géométrie et propriétés physiques. \\
\bottomrule
\end{tabular}
\par\vspace{.22cm}
{\footnotesize L'apport propre d'Hypergraph-Percol reste à mesurer.}
\reffoot{spt23,bps19,poly25}
\end{frame}
